\chapter{Chemical Filtering of Candidate Polymers}
\label{sec:filtering_chapter}

Complementing the quantitative predictions of the MLP, a set of heuristic filters based on established chemical principles is applied to eliminate polymers that are unlikely to be effective adsorbents. These filters encode domain knowledge about the requirements for successful adsorption in wastewater treatment.

\section{Filter Criteria}

The first filter concerns the partition coefficient (logP) of the polymer. For a polymer to function as an insoluble adsorbent in aqueous solution, it must have a sufficiently high logP value to remain in a solid phase. A minimum logP threshold of 1.5 is applied, ensuring that the polymer remains hydrophobic enough to avoid dissolution while not being so lipophilic that it loses affinity for the aqueous pollutant.

The second filter addresses polar surface area (TPSA). Polymers with higher TPSA values possess more polar functional groups that can engage in hydrogen bonding and electrostatic interactions with pollutant molecules. A minimum TPSA of 60 \r{A}² is required to ensure sufficient polar character for effective binding.

The third and most important filter is based on Frontier Molecular Orbital (FMO) theory. The HOMO and LUMO energies of both the polymer and the target drug determine the nature of their electronic interactions. Two interaction pathways are possible: the polymer acts as an electron donor (using its HOMO) to the drug (acting as an acceptor via its LUMO), or vice versa. The energy gaps for these two pathways are computed as:

\begin{equation}
	\Delta E_1 = |E_{\text{LUMO, drug}} - E_{\text{HOMO, polymer}}|
\end{equation}

\begin{equation}
	\Delta E_2 = |E_{\text{LUMO, polymer}} - E_{\text{HOMO, drug}}|
\end{equation}

The smaller of these two gaps represents the dominant interaction pathway. A smaller intermolecular FMO gap indicates stronger electronic coupling and thus stronger binding. A maximum gap threshold of 4.0 eV is applied, retaining only those polymer-molecule pairs with favorable electronic complementarity.

The fourth filter addresses synthetic accessibility (SA score). A practical wastewater adsorbent must be synthesizable at scale using reasonable chemical methods. The SA score, computed from structural features that affect synthetic difficulty, is constrained to a maximum value of 4.5 on a 1-10 scale.

These filters are applied sequentially after generation and validation, dramatically reducing the number of candidate polymers before they are passed to the predictive model for capacity estimation. The combination of chemical filtering and machine learning prediction provides a balanced approach that leverages both domain knowledge and data-driven pattern recognition.
